% 03-Apparatus.tex <- v1 sec:methods-data. Numbers verbatim from v1.
% NO claim block: the absence is structural.

\chapter{The atlas, the model and the rules of measurement}
\label{sec:methods-data}
\framedecl{ER}

\begin{question}
What has to be settled before a score on this atlas means anything? Four
things. What is measured, on what population, against what reference, with what
uncertainty. Tahoe answers none by convention, so each is settled below by a
measurement, and each narrows a claim later.
\end{question}

\section{The panel is a budget, not a biology}

Tahoe-100M~\cite{tahoe100m} holds drug-treated single-cell transcriptomes across
cancer cell lines, with plate-matched DMSO controls. Metadata supplies cell line
(Cellosaurus ID), drug, dose (\texttt{drugname\_drugconc}), plate, and a
fine-grained mechanism-of-action annotation (\texttt{moa-fine}), six fields that
are the whole of what any prompt below can carry.

All analysis is restricted to the $G = 946$ genes in the intersection of the
L1000 landmark set~\cite{l1000} with Tahoe's measured genes, library-normalised
to counts per $10{,}000$ and transformed as $\log_{10}(1 + x)$. The restriction
is a budget choice, not a biological one: the context is $8192$ tokens and a
prompt carries an instruction line and two cell sentences, so the transcriptome
does not fit. Normalisation multiplies every gene in a cell by one constant and
the logarithm is monotone, so within a cell the ordering, and therefore the
sentence, is invariant to both.

\paragraph{What it does change} Everything downstream of a single cell. Scaling
by library size improves comparability across cells of different depth,
but the profiles remain compositional, not absolute, so pseudobulk means, the
drug-specific residual, and every cosine and Euclidean quantity built on them
depend on the normalisation even though the sentences do not. Depth also shapes
which genes are detected, though never their order.

\section{Only abundance checks gene identity}

Tahoe's per-cell \texttt{genes} field holds vocabulary token identifiers, not
positional indices into the gene metadata table. Mapping them positionally
scrambles gene identity quietly; the values exceed the metadata row count, so
the error surfaces only as an out-of-bounds index at the tail. All expression
construction here maps \texttt{gene\_metadata['token\_id']} to gene symbol to
panel index.

The obvious check cannot fail. Permuting the panel's columns is an orthogonal
transform of the row space, so distances, principal components and accuracy are
preserved exactly, and a PCA cell-line probe returns $0.864$ against a chance
rate of $0.02$ on the correct panel and $0.864$ on a scrambled one. What
discriminates is abundance, a property of the gene and not of the panel.
\texttt{GAPDH} is the most abundant of the $946$ genes and \texttt{RPS6} the
second, and across the canonical high-expressors the median abundance rank is
$19.5$, against $473.5$ expected under a permuted index ($p = 0.00015$,
$20{,}000$ draws; \texttt{gene\_identity\_check.py}). The probe that survives
the permutation is the probe that cannot detect it.

\section{Every treatment sits in one well}

Tahoe applies one drug at one dose to one \emph{well} holding roughly $49$ cell
lines. This project's cache spans more plates than the three its scored
conditions fall on (\texttt{plate4}, \texttt{plate5}, \texttt{plate6}), and
holds $289$ treated wells, a mean of $48.6$ cell lines each (median $49$, range
$43$--$50$), with $286$ of $287$ distinct (drug, dose) treatments occurring in
exactly one well.

Those figures describe this cache, not the atlas: Tahoe includes further plates,
among them a replication plate, so ``one well'' below means one well
\emph{here}. The single exception, one treatment spread over three wells, is
also the whole of the cross-plate structure at condition level, and it accounts
for the $49$ (drug, dose, cell line) groups the unit audit finds on more than
one plate. Every other condition on more than one plate differs in dose, and
each drug is assayed on its own plate or plates with plate-matched controls.

\paragraph{Three consequences} Conditions from one well share a physical
assignment, so they are not independent (\cref{sec:apparatus-clusters}); drug
and plate are partially confounded, so discrimination is measured within plate
(\cref{sec:res-blind}); and a (drug, cell line) pair cannot be held out cleanly
(\cref{sec:methods-holdout}). \cref{fig:confound} draws the three in that
order.

\begin{figure}[htbp]
\centering
% ===========================================================================
% figs/confound.tex -- fig:confound, "One assignment, many conditions".
% Figure BODY ONLY. \input this inside a figure environment; the caption and
% \label stay in Sections/03-Apparatus.tex.
%
% SCHEMATIC. Nothing here is plotted from data. The three quantities it names
% are the ones the surrounding text states: roughly 49 cell lines per treated
% well, the other roughly 48 that stay in training when one pair is held out,
% and plate4/plate5/plate6. The squares are twelve per well for legibility,
% which the caption says out loud so the count is not read as a datum.
%
% ---- METRIC REWORK, and why the numbers below are what they are ------------
% The v1 drawing measured 354.51pt (124.6mm) inside a 404.03pt (142mm) block:
% it under-filled the measure by 12%. Fixing that is NOT a uniform \scale of
% the picture, and it is not the 1.44 a naive width ratio suggests, because
% the picture is half drawing and half fixed-size type. Only the drawing
% scales; the three right-hand glosses and the (a)/(b)/(c) labels are set in
% \scriptsize and must stay at that size, since they already print at the
% document's apparatus size.
%
% Measured with pdflatex over this preamble (\the\wd of a savebox):
%   widest gloss node  "a cross-plate drug ... plate contrast" = 181.46pt
%   "held out" node half-width                                 =  18.75pt
% With k the scale on the drawn geometry, the bounding box spans
%   minx = 0.42k*28.4527 - 18.75      (the "held out" label overhangs left)
%   maxx = (5.65k + 0.4)*28.4527 + 181.46
%   width = 148.807k + 211.596 pt
% so k = 1.286 lands the picture at 402.9pt. VERIFIED at 402.96pt; the block
% is 404.03pt. A pure similarity transform of the drawing: x and y take the
% same unit so the lanes keep their aspect, and the squares' minimum size is
% scaled with them (3.0mm * 1.286 = 3.86mm) so a widened lane does not end up
% holding a sparse row of undersized marks. Coordinate VALUES are untouched --
% every number below is v1's -- which is what keeps this a metric change and
% not a redraw.
%
% Gloss alignment: all three glosses sit at one x, the rightmost lane's right
% edge (5.65 units) plus a physical 4mm. 4mm is 0.4cm/1.286 = 0.311 units, so
% x = 5.961. Strip (b)'s own right edge is 5.50, short of (a) and (c) at 5.65;
% hanging its gloss off its own edge would step the left margin of the three
% glosses in and out for no reason.
%
% NEVER \resizebox this picture. Changing its printed size means changing k
% here and re-measuring, because the type must not move.
% ===========================================================================
\begin{tikzpicture}[x=1.286cm, y=1.286cm, font=\small,
  lane/.style={draw=rulegrey, line width=0.4pt, rounded corners=1pt},
  cellb/.style={draw=rulegrey, line width=0.3pt, fill=panelbg,
                minimum size=3.86mm, inner sep=0pt},
  outb/.style={draw=accent, line width=0.7pt, fill=white,
               minimum size=3.86mm, inner sep=0pt},
  lbl/.style={font=\scriptsize, text=quietink},
  strong/.style={font=\scriptsize\bfseries, text=ink}]

% (a)
\node[strong, anchor=west] at (0,3.95)
  {(a)\quad one drug at one dose enters one well};
\draw[lane] (0.15,3.10) rectangle (5.65,3.60);
\foreach \i in {0,...,11}{\node[cellb] at (0.42+\i*0.44,3.35) {};}
\node[lbl, anchor=west] at (5.961,3.35)
  {$\approx 49$ conditions, one assignment};

% (b)
\node[strong, anchor=west] at (0,2.55)
  {(b)\quad each drug is assayed on its own plate};
\foreach \p/\x in {4/0.15, 5/2.00, 6/3.85}{
  \draw[lane] (\x,1.60) rectangle (\x+1.65,2.10);
  \node[lbl] at (\x+0.825,1.85) {\texttt{plate\p}};}
\node[lbl, anchor=west] at (5.961,1.85)
  {a cross-plate drug contrast is also a plate contrast};

% (c)
\node[strong, anchor=west] at (0,1.05)
  {(c)\quad a held-out pair leaves its well in training};
\draw[lane] (0.15,0.35) rectangle (5.65,0.85);
\node[outb] at (0.42,0.60) {};
\foreach \i in {1,...,11}{\node[cellb] at (0.42+\i*0.44,0.60) {};}
\draw[accent, line width=0.4pt] (0.42,0.32) -- (0.42,0.10);
\node[lbl, anchor=north, text=accent] at (0.42,0.08) {held out};
\node[lbl, anchor=west] at (5.961,0.60)
  {its other $\approx 48$ conditions stay in};

\end{tikzpicture}%

\caption[The well is the unit the atlas actually assigns]{\textbf{One
assignment, many conditions.} The roughly $49$ conditions from one well share
an assignment and a batch (a). Each drug is assayed on its own plate or plates,
so a drug contrast across plates is also a plate contrast (b). Removing a (drug,
cell line) pair leaves its well in training through the other $\approx 48$ lines
(c). Squares are conditions, twelve per well for legibility.}
\label{fig:confound}
\end{figure}

\begin{defn}{Condition, treated well, batch unit}
A \emph{condition} is a (drug, cell line, plate, dose) tuple, and it is the unit
everything below is scored at. The physical assignment is the \emph{treated
well}, so a condition sits inside an assignment and is not one. Plates are
numbered within cell line, not globally, so the batch unit is the pair (cell
line, plate), never the plate label alone. A \emph{pseudobulk} is a group's mean
expression profile; a \emph{control pseudobulk} is the same over plate-matched
DMSO cells.
\end{defn}

Dose is a molar concentration resolved from \texttt{drugname\_drugconc}; an
earlier pipeline kept only the numeric part, which collides $1\,\mu$M with
$1\,$nM. An audit of all $289$ treated wells found $289$ usable molar doses
with zero unit collisions, so the defect was a demonstrated possibility that did
not materialise. Combination treatments are dropped, not analysed as their first
component.

\begin{scopelimit}[carried into \cref{sec:methods-holdout} and \cref{sec:limitations}]
A (drug, cell line) pair cannot be held out cleanly on this cache. Removing the
pair removes one condition, and its well remains in training through the other
$\approx 48$ lines. The residual arm therefore builds its own, larger holdout.
\end{scopelimit}

\section{Barely half the atlas is identifiable}

\cref{tab:scale} states the quantities every $n$ in the investigation should be
located against. One symbol in it runs ahead of its
derivation.\seealso{\cref{sec:res-metrics} derives $\NIR$ and audits its chance
level with empirical control arms.} $\NIR$, the normalised inverse rank, is the
fraction of a declared comparison set whose truth a prediction resembles less
than its own, so chance is $0.50$ under the exchangeability condition of
\cref{sec:res-metrics}.

\begin{scopelimit}[carried into \cref{sec:res-blind} and \cref{sec:res-encoding}]
Two populations are used and they are not interchangeable. The broad
characterisation (\cref{sec:res-blind}) streams a wide sample of the atlas, over
more plates and drugs than cached, so the drug and condition counts in its block
of \cref{tab:scale} count that sample, not the cache. The residual arm
(\cref{sec:res-encoding} onward) works from a smaller cached subset. A count
from one block does not bound a count in the other.
\end{scopelimit}

\begin{table}[htbp]
\begin{fullwidth}
\begin{threeparttable}
\caption{\textbf{The populations behind every $n$ quoted in the investigation.}
The upper block is a streamed sample of the atlas; the lower two are the cache
built once and reused.}
\label{tab:scale}
\small
\centering
\begin{tabular}{p{0.64\linewidth}r}
\toprule
\multicolumn{2}{l}{\thead{Task characterisation} (streamed sample)} \\
\midrule
conditions characterised & \num{6628} \\
\hangindent1em\hangafter0 of which statistically active against DMSO & 3,647 / 5,014 (72.7\%) \\
\hangindent1em\hangafter0 empirically identifiable at the tested pseudobulk budget (own within-plate
  split-half $\NIR \geq 0.8$) & 3,064 (46.2\%) \\
\hangindent1em\hangafter0 with a plate-mate closer than their own split-half & 78.7\% \\
median cells per condition & 44 \\
median \#DEGs per condition (BH $q<0.05$) & 0 \\
median SNR (effect $\div$ within-well split-half noise) & 0.75 \\
drugs with a real dose series & 218 / 351 (62.1\%) \\
drugs seen in $\geq 3$ cell lines & 227 \\
\midrule
\multicolumn{2}{l}{\thead{Residual-frame cache} (rebuilt targets,
  \cref{sec:methods-residual})} \\
\midrule
cells (treated / control) & 620,564 (600,000 / 20,564) \\
(cell line, plate) groups behind scored conditions & 119 \\
drugs / cell lines & 106 / 50 \\
conditions with $\geq40$ treated and $\geq20$ control cells & 6,617 \\
\hangindent1em\hangafter0 by split after inventory filters (train / held-out wells / held-out drugs)
  & 4,999 / 900 / 711 \\
\hangindent1em\hangafter0 retained by the data-resolved reliability line & 2,523 of 4,999 train
  (50.5\%) \\
training examples / well-held-out validation examples & 147,710 / 3,600 \\
\midrule
\multicolumn{2}{l}{\thead{Evaluation sets}} \\
\midrule
tier 1 (seen conditions) / tier 2 (unseen drugs) / tier 3 (unseen combinations)
  & 4,325 / 2,188 / 99 \\
conditions scored per tier, generation evaluations & 400 \\
\raggedright tier-2 conditions in the expression-frame $\NIR$ benchmark
  (within plate) & 606 \\
conditions scored in the residual-frame evaluation & 1,394 \\
\hangindent1em\hangafter0 by split (train / \texttt{unseen\_combo} / \texttt{unseen\_drug})
  & 300 / \needsaudit{895} / 199 \\
\hangindent1em\hangafter0 treated wells behind them (no well spans two splits) & 136 / 36 / 28 \\
\hangindent1em\hangafter0 drugs per split (90 distinct in total) & 76 / 34 / 10 \\
\hangindent1em\hangafter0 cell lines per split (42 distinct in total) & 39 / 42 / 40 \\
\bottomrule
\end{tabular}
\begin{tablenotes}[flushleft]\footnotesize
\item The $72.7\%$ figure is over the $5{,}014$ conditions whose (cell line,
plate) group carries plate-matched DMSO cells, not all $6{,}628$.
\item The identifiability bar is per-condition, distinct from the aggregate
ceilings of \cref{tab:ceilings}.
\item The \texttt{unseen\_combo} count is flagged pending a cluster rerun; no
argument turns on its exact value.
\end{tablenotes}
\end{threeparttable}
\end{fullwidth}
\end{table}

\paragraph{The upper block is a result} It describes Tahoe-100M, not this
project. Across the $5{,}014$ control-reachable conditions the median has no
differentially expressed gene at BH $q<0.05$ and a signal-to-noise ratio of
$0.75$ against its own within-well split-half noise, while $72.7\%$ are active
on the whole profile, leaving $27.3\%$ statistically inert. A per-gene test is
far harsher, and $27.3\%$ alone would make the atlas sound better resolved than
it is. The typical condition carries a real response
measured with too few cells, and not no response (\cref{sec:res-blind}).

The identifiability bar is per-condition: a condition's own split-half $\NIR$
against its plate-mates has to clear $0.8$. Of the $6{,}628$ conditions,
$46.2\%$ clear it, and $78.7\%$ have a plate-mate lying closer to them than
their own second half does. Identifiability is a property of the pairing, not of
the drug: over the $227$ drugs assayed in at least three cell lines, the median
drug's split-half reference spans $0.83$ between its most and least identifiable
context (\texttt{drug\_biology\_atlas.json}). A benchmark omitting those figures
quotes accuracies over a population barely half of which is empirically
identifiable at the tested pseudobulk budget, and a per-drug leaderboard on it
ranks pairings.

% The term is first used in tab:scale above, inside a float, where a marginpar
% cannot be set. This is its first prose use and the chapter's definition point,
% per the \gloss convention (definition at first use in each chapter). Wording
% is the one 05-Limitations.tex already uses, so the two agree word for word.
\paragraph{A filter this table does not total} The middle block records
$2{,}523$ of $4{,}999$ training conditions as \emph{retained by the
data-resolved reliability line},\gloss{reliability line}{The half-split cosine
threshold above which a condition's residual counts as reproducible.} which
\cref{sec:methods-residual} resolves from the data and applies to the training
split alone. It does not state the pool
once the held-out splits join that filtered set. The channel gate of
\cref{sec:res-scope} later scores \needsaudit[the retained pool]{4{,}131}
retained conditions; a cluster rerun settles whether that is this table's
three-way tally under another name, and no figure here depends on it.

The three tiers of the last block are defined by what training contained.
\emph{Tier 1} holds conditions whose (drug, cell line) pairing was
seen in training, so it measures fit and not generalisation. \emph{Tier 2} holds
conditions whose drug never appeared, Leave-Drugs-Out in the vocabulary of
\cref{sec:lit-dreval}. \emph{Tier 3} holds pairings whose drug and cell line
were each seen but never together, Leave-Pairs-Out; it is small ($99$
conditions) and is reported only where adequately powered. The residual arm
builds its own, larger holdout instead: \cref{sec:apparatus-vocab} names its
three splits, and \cref{sec:methods-holdout} settles the unit they are keyed
at.

\section{The backbone has seen no perturbation}

A cell is written as a \emph{cell sentence}, its expressed panel genes in
descending order, terminated by \ENDCELL. The prompt supplies cell line, drug,
dose, mechanism and the control cell's sentence; what the model must produce is
the treated cell's. \ENDCELL{} and (for the residual arm) \DOWN{} are registered
as atomic special tokens; unregistered, \DOWN{} splits into sub-words and the
up/down boundary stops being a clean symbol.

The backbone is \texttt{C2S-Scale-Pythia-1b-pt}~\cite{c2sscale}, a C2S-Scale
checkpoint on the Pythia-1b decoder~\cite{pythia}, cold-started for every arm
from the same base checkpoint with identical hyperparameters (1 epoch, learning
rate $10^{-5}$, gradient accumulation 16, maximum length 8192,
bf16).\seealso{\cref{sec:appendix} tabulates the training arms.} Cold-starting
makes the arms comparable; the training target or objective is the only
variable.

\paragraph{What the checkpoint has seen} It is pretrained on over 50 million
human and mouse cells from 825 single-cell datasets in CELLxGENE and the Human
Cell Atlas, an atlas corpus with no chemical or genetic perturbation data. Its
model card lists cell and tissue prediction, conditioned generation and gene-set
conversion, but not perturbation prediction. The base model arrives at
fine-tuning here having never seen a drug-treated cell.

\paragraph{Two channels stay open} The decoder underneath is Pythia-1b,
pretrained on text in which drug names, targets and mechanisms occur, and every
prompt supplies the drug's \texttt{moa-fine} annotation. A drug the fine-tuning
never saw is not a drug the system knows nothing about.

\begin{scopelimit}[carried into \cref{sec:res-scope} and \cref{sec:limitations}]
``Held out'' therefore means held out from Tahoe fine-tuning, with mechanism
supplied, and every unseen-drug result below is read under that narrower scope.
\end{scopelimit}

\section{The generic replaces the target}
\label{sec:apparatus-vocab}

Two \emph{measurement frames} run through the investigation and are never
interchangeable (\cref{fig:frames}): the \emph{expression frame} scores a
predicted cell profile against other drugs on the same plate, the
\emph{residual frame} (\cref{sec:methods-residual}) scores the drug-specific
residual against other drugs in the same cell line. Both put chance at $0.50$. The
\emph{comparison set} $J$, the conditions a prediction is ranked against, is
what fixes it there, and it is declared at every score below. A score without
its comparison set is not a score.

% ===========================================================================
% frames.tex -- the two measurement frames, EXPR and RESID, as two cards.
% \input by Sections/03-Apparatus.tex. Compiles against thesis_v2/preamble.tex
% and nothing else; no external data, no \includegraphics.
%
% NO DATA SOURCE. This is a schematic. The only numeral on the cards is the
% chance level 0.50, which is a definition and not a measurement: both frames
% are two-alternative retrievals, so chance is 0.50 by construction. It agrees
% with the prose immediately above this figure in Sections/03-Apparatus.tex
% ("Both put chance at $0.50$"), with Sections/02-Background.tex:244 and with
% Sections/04a-ruler.tex:530 ("Chance is $0.50$ by ..."). Nothing here is read
% from RESULTS_cluster/ because nothing here is a result. The numeric version
% of this argument -- the two frames on their own scales -- is fig:two-scales
% in 04d; Chapter 3 is apparatus and comes before those results exist, so this
% figure carries no scale, no axis and no measured quantity. Sibling
% frames.json records the geometry and that single number.
%
% ---------------------------------------------------------------------------
% THE BUG, AND WHY THE FIX IS NOT THE ONE-LINER IT LOOKS LIKE.
%
% Symptom, both cards: the two-line value on row 2 collided with its own key.
% Measured in the built thesis, ink gap between the key and its value, in
% PostScript points: r1 -0.02, r2 -3.02, r3 +2.12. Only r2 is visible as a
% collision. (r1 is the descender of "predicted" level with the ascender of
% "profile"; different x, they do not touch.)
%
% Diagnosis, as briefed and confirmed: val/.style used anchor=WEST, which
% centres a node's box on its y, so a TWO-LINE value grew upward into the key
% above it while one-line values did not.
%
% The briefed fix was anchor=north west with every value at its key's
% y - 0.30cm. Built, that removes the collision and introduces a worse fault:
% it lowers each value until it sits CLOSER TO THE NEXT KEY THAN TO ITS OWN,
% so every value reads as the heading of the row beneath it. Sweeping the
% offset, all measured inside one PDF by tools/_measure_frames.py and recorded
% in frames.json, gave "within" = key to its own value, "between" = value to
% the next key:
%
%   offset      within r1 / r2 / r3      between: after r1 / after r2
%   0.15          2.21  2.24  3.14            9.60        7.43
%   0.30          6.46  6.49  7.39            5.35        3.18   <- reversed
%   0.40          9.29  9.33 10.23            2.51        0.35
%
% and then the same figure BUILT INSIDE THE THESIS at offset 0.15 measured
% 5.95 / 5.98 / 5.29 within against 5.86 / 3.69 between -- 3.7pt away from the
% harness, and ambiguous. The offset is not a constant of the figure.
%
% ROOT CAUSE, and why this file no longer uses align=left. TikZ's `align'
% option typesets the node in a minipage, and a minipage's first-line height
% follows the AMBIENT \baselineskip. So a value node's ink lands in a place
% that depends on the prose that happens to precede the float: about 2.2pt
% between a \normalsize paragraph and a \small line, and about 3.7pt between a
% test harness and Chapter 3. No choice of offset survives an edit to the
% surrounding text. Calibrating one would have hidden the fault, not fixed it.
%
% THE FIX ACTUALLY USED: every node is a single line anchored BASE WEST, so
% each y coordinate below IS that line's baseline. Nothing is centred, nothing
% is measured from a box edge, no minipage is created, and no ascender or
% descender can move a line. The two-line value is two nodes rather than one
% node with `\\'. Key-to-value leading is therefore identical on all three rows
% and on both cards by construction -- which is what the brief asked for -- and
% it is now also identical from build to build.
%
%   key baselines    3.70   2.72   1.52     (row pitch unchanged)
%   value baselines  3.37   2.39   1.19     (key - 0.33cm)
%   row-2 value, second line  2.06          (previous line - 0.33cm)
%
% ONE step governs everything inside a row: 0.33cm = 9.39pt, one \scriptsize
% line of leading. A key and its value therefore set as two consecutive lines
% of a single block, exactly as rows 1 and 3 of the published figure did, and a
% wrapped value continues on the same rhythm. Row SEPARATION is carried
% entirely by the key pitch, which is untouched (0.98cm, then 1.20cm), giving
% 0.65cm and 0.53cm of clear space before the next key -- roughly twice the
% within-row step, so a value binds to the key above it and not to the one
% below.
%
% Verified by reading the baselines back out of the built PDFs. Chapter 3 in
% main.pdf and all four contexts of figs/_frames_harness.tex report the same
% seven baselines to 0.001cm:
%   3.703  3.375 | 2.727  2.398  2.059 | 1.531  1.203   (tikz y, both cards)
% Note the harness and the thesis are only comparable on BASELINES. Glyph
% bounding boxes reported by a PDF text extractor differ by a constant ~2.1pt
% between the two files because the font subsets differ, so ink-gap numbers
% are only meaningful WITHIN one PDF.
%
% OTHER LAYOUT NOTES
%   * cards are 5.9cm wide with a 0.5cm gutter, 12.3cm overall. NOT rescaled:
%     12.3cm sits inside the 14.2cm body block and leaves the page free to
%     carry a margin note.
%   * the dotted rule and the sentence under it span BOTH cards on purpose.
%     They are the only mark that crosses the gutter; everything else is
%     card-local, because the point of the figure is that the columns do not
%     exchange.
%   * if a value ever needs a third line, add a node at (its predecessor
%     - 0.34cm). Do not reintroduce align/`\\'.
% ===========================================================================
\begin{figure}[htbp]
\centering
\begin{tikzpicture}[x=1cm, y=1cm, font=\small,
  panel/.style={draw=rulegrey, line width=0.4pt, rounded corners=1pt},
  % base west on both: every y below is a BASELINE. See the header note --
  % anchor=west centres the box (that was the collision) and align=left makes
  % the node's height depend on the surrounding prose (that was worse).
  key/.style={font=\scriptsize\bfseries, text=quietink, anchor=base west},
  val/.style={font=\scriptsize, text=ink, anchor=base west}]

\draw[panel] (0,0.55) rectangle (5.9,4.6);
\draw[panel] (6.4,0.55) rectangle (12.3,4.6);

\node[anchor=base west] at (0.25,4.18) {\Eframe};
\node[anchor=base west] at (6.65,4.18) {\Rframe};

% key baseline 3.70 / 2.72 / 1.52 ; value baseline = key - 0.33 ; wrap - 0.34.
\node[key] at (0.25,3.70) {what is predicted};
\node[val] at (0.25,3.37) {a cell profile};
\node[key] at (6.65,3.70) {what is predicted};
\node[val] at (6.65,3.37) {the drug-specific residual};

\node[key] at (0.25,2.72) {what it is ranked against};
\node[val] at (0.25,2.39) {other drugs on the};
\node[val] at (0.25,2.06) {same plate};
\node[key] at (6.65,2.72) {what it is ranked against};
\node[val] at (6.65,2.39) {other drugs in the};
\node[val] at (6.65,2.06) {same cell line};

\node[key] at (0.25,1.52) {where chance sits};
\node[val] at (0.25,1.19) {$0.50$};
\node[key] at (6.65,1.52) {where chance sits};
\node[val] at (6.65,1.19) {$0.50$};

\draw[ink, line width=0.4pt, densely dotted] (0,0.25) -- (12.3,0.25);
\node[font=\scriptsize\itshape, text=ink, anchor=base] at (6.15,0.02)
  {the same number in the two columns is not the same measurement};
\end{tikzpicture}
\caption[The two measurement frames]{\textbf{Two frames, one chance level, and
no exchange rate between them.} The prediction, the comparison set and the
meaning of a score change between frames; the chance level does not.}
\label{fig:frames}
\end{figure}


\begin{defn}{The generic and the residual target}
The \emph{generic} is the drug-agnostic mean response subtracted to form the
residual, operationally the mean shift over the other drugs at the stated scope,
conventionally read as the stress and cell-cycle programme common to all
drugs --- a reading used for intuition and nowhere tested. The
\emph{residual target} is what remains. It is built \emph{split-before-fit}
(holdout fixed from metadata before any statistic is estimated) and
\emph{leave-one-drug-out} (no condition's generic contains its own drug). The \emph{scope}, \texttt{cell\_line} or \texttt{plate}, is the set
of drugs the mean is taken over.
\end{defn}

The generic carries $39.0\%$ of the perturbation shift's energy
(\cref{sec:res-transfer}), so subtracting it does not clean the target up but
replaces it: every residual-frame number is a score on what is left.

\paragraph{Baselines and splits} \emph{Control-copy} predicts the untreated cell
unchanged; carrying zero drug information, it is the floor a drug-aware model
has to clear. The residual arm's split has three parts,
\texttt{train}, \texttt{unseen\_combo} (an unseen \emph{well} of a seen drug)
and \texttt{unseen\_drug}, sampled under an explicit per-split \emph{quota},
because a uniform draw reproduces the pool's proportions and starves the
held-out splits (\cref{sec:methods-holdout}). Five training \emph{arms} differ
only in target or objective; where a claim is multiplicity corrected, the family
is those five by up to seven depths.

\section{One estimator, and not a replicate}

Almost every negative result below is stated as a distance from a
\emph{ceiling}. The word names one estimator, not seven.

\begin{defn}{Within-well split-half precision reference}
The score obtained when a disjoint cell subset of the same treated well is
substituted for the prediction, recomputed for every stratum scored, so its
value moves with the stratum. It is not a replicate ceiling. What varies across
the rows of \cref{tab:ceilings} is the frame, the comparison set and the cell
budget, never the estimator.
\end{defn}

Two properties limit how the ceilings read. First, they are \emph{split-half}
estimates and not biological replicates, because this cache places $286$ of
$287$ treatments in exactly one well. Second, a split-half ceiling estimates
within-well sampling precision, which is optimistic relative to true biological
reproducibility, and every coverage figure in \cref{sec:res-lookup} inherits
that optimism.

\paragraph{The one exception} One ceiling is not that estimator and not a row in
\cref{tab:ceilings}. The positive control inside the expression-frame $\DRF$
audit of \cref{sec:res-metrics} is a denoised interpolated duplicate, not a raw
split-half, and that section turns on the difference.

\begin{table}[htbp]
\begin{fullwidth}
\begin{threeparttable}
\caption{\textbf{One estimator, seven settings.} Every ``ceiling'' here is a
within-well split-half precision reference, the score of a disjoint subset of
the same treated well's cells. They differ in frame, comparison set, cell budget
and split contents.}
\label{tab:ceilings}
\footnotesize
\begin{tabularx}{\linewidth}{@{}>{\raggedright\arraybackslash}p{0.18\linewidth}
  >{\raggedright\arraybackslash}p{0.20\linewidth}
  >{\raggedright\arraybackslash}p{0.31\linewidth}X@{}}
\toprule
\thead{Ceiling} & \thead{Frame and metric} & \thead{Comparison set and budget}
  & \thead{Where used} \\
\midrule
$0.76$ & expression, $\DEdr$ & two real cells & metric audit
  (\cref{sec:res-metrics}) \\
$0.625$ / $0.575$ & expression, $\NIR$ & leak audit; cross-plate / within plate
  & plate leak (\cref{sec:res-blind}) \\
$0.67$--$0.83$ & expression, forced choice & within plate, by cell budget
  & \cref{sec:res-blind} \\
$0.576$ & expression, $\NIR$ & within plate, tier 2 & the blindness verdict \\
$0.61$ / $0.76$ / $0.85$ & expression, $\NIR$ & within plate;
  $10$/$40$/$120$ cells & identifiability curve \\
$0.956$/$0.824$/$0.836$ & residual, $\NIR$ & within cell line;
  train/combo/drug & \cref{sec:res-lookup} \\
$0.742$ / $0.7432$ & token retrieval & within condition, split-half
  & \cref{sec:methods-residual} \\
\bottomrule
\end{tabularx}
\begin{tablenotes}[flushleft]\footnotesize
\item The $\DRF$ positive control of \cref{sec:res-metrics} is the
investigation's one exception and is not a row here.
\item None of these is a biological replicate ($286$ of $287$ treatments sit in
a single well).
\item In the residual frame the raw reference is not constant across splits (a
spread of $0.132$), but that spread is an artefact of how the splits were
filtered and not a difference in recoverable signal; matched at the training
reliability line the three agree to $0.006$ (\cref{sec:limitations}).
\end{tablenotes}
\end{threeparttable}
\end{fullwidth}
\end{table}

\begin{scopelimit}[carried into \cref{sec:res-lookup} and \cref{sec:limitations}]
A distance to the reference is comparable within a split and not between them,
because the splits are scored on populations selected by different rules. And
because the reference is a sampling-precision estimate and not a biological
replicate, every coverage figure computed against it is a distance to an
optimistic bar by an unknown amount. Where one is quoted, the ordering carries
the claim.
\end{scopelimit}

\section{Two crossed dependencies, not one}
\label{sec:apparatus-clusters}

Two rules apply to every result below, plus one statement of provisionality.
Conditions are not independent, and not independent in two crossed ways at once
(\cref{fig:clustering}): two from one treated well share a physical assignment
and a batch, two from one cell line share the biology. Treating them as
independent gives intervals about $1.5$ times too narrow on the scramble
strata and $1.7$ to $2.5$ times too narrow on the contrasts against the lookup
baselines. Clustering on the assignment unit alone uses the larger cluster
count ($289$ wells against $\approx 50$ lines) and would tighten every interval
while sounding more rigorous. That is the failure the estimator below avoids.

% ===========================================================================
% clustering.tex -- the two crossed groupings, plus the drug the wells nest in.
% \input by Sections/03-Apparatus.tex. Compiles against thesis_v2/preamble.tex
% and nothing else; SCHEMATIC -- no external data, no \includegraphics, and
% deliberately NO number anywhere in the drawing. The lattice is 9 columns of
% 6 dots because that is what reads at this measure, not because the cache
% holds 9 wells or 6 lines; the caption says so in as many words.
%
% WHY IT WAS REDRAWN. The version it replaces sat inline in 03-Apparatus.tex
% at lines 477-510 (v1's is under ../thesis/ and is not touched).
%
%   1. THE DOT WAS NOT ON THE LATTICE. The old bands were centred at y=2.6 and
%      x=5.0, while the lattice rows sat at 0.5+0.6r (2.3, 2.9) and the columns
%      at 0.7+1.1c (4.0, 5.1). Both centre lines fell BETWEEN dots, so the
%      accent "one condition" dot was drawn in the gap. The caption said "the
%      dot where they meet is one condition" and the drawing did not show one.
%      Here every coordinate is derived from the two constants the lattice is
%      generated from (\px, \py) and the highlight is addressed by INDEX
%      (\hc, \hr), so band centres and accent dot cannot drift off a lattice
%      point again.
%
%   2. THE BANDS WERE PADDED BY EYE. Old horizontal band 0.35..8.45 against
%      dots at 0.7..7.3 -- 0.35 short at the left, 1.15 long at the right. Old
%      vertical band 0.12..4.08 against dots at 0.5..4.1 -- 0.38 short at the
%      bottom, 0.02 over at the top. Here each band ends exactly half a lattice
%      pitch past its first and last dot (\BL,\BR,\BB,\BT below), computed, so
%      both are symmetric by construction.
%
%   3. THE DOTS WERE TEXTURE, NOT CONDITIONS. They were rulegrey!70, i.e.
%      #9A9A9A at 70% over white ~= #C0C0C0 -- lighter than the document's
%      quiet ink and far too light for a mark the caption asks the reader to
%      count. They are quietink (#6B6B6B) here, and 1.7pt rather than 1.5pt.
%
%   4. IT PAID FOR FULLWIDTH AND DID NOT USE IT. The old float opened
%      \begin{fullwidth} -- 174mm, and by the preamble's own rule a page
%      carrying one may then carry no margin note -- and drew 148.6mm. This
%      draws inside the 142mm measure (see MEASURED below), so it is an
%      ordinary body-width float and its page keeps the margin column.
%
%   5. THE DRUG WAS IN NO FIGURE AT ALL, and it is the axis that flips three
%      verdicts later. It is drawn in a deliberately different register from
%      the two bands: outline only, no fill, rounded corners, against two solid
%      tint rectangles. A well carries one drug at one dose (sec. "Every
%      treatment sits in one well"), so wells NEST inside drugs -- three whole
%      well-columns per envelope -- while the cell-line band CROSSES all three
%      envelopes and runs past the outer edge of both end ones. Nesting and
%      crossing are then the same picture, told apart by shape, not by colour.
%
% GEOMETRY, in cm, all of it derived from \px and \py. Every clearance here is
% computed from those two constants, not judged by eye:
%   lattice     9 cols x 6 rows, pitch 1.25 x 0.81, dots (0,0)..(10.00,4.05)
%   highlight   col 4 -> x = 5.00, row 4 -> y = 3.24     <- both ON the lattice
%   bands       half-thickness 0.25, so 0.155 of ground to the next dot row
%               (0.405 half-pitch - 0.25) and 0.095 clear of that dot's edge;
%               ends at -0.625/10.625 and -0.405/4.455, i.e. exactly half a
%               pitch past the first and last dot of the row and the column
%   envelopes   3 x 3 columns, pad 0.46 -> -0.46..2.96, 3.29..6.71, 7.04..10.46,
%               so 0.33 between neighbours; vertically -0.585..4.635, giving
%               0.18 of margin above and below the well band each one contains
%   callout     the leader leaves the accent dot at 46.3 deg, clears the
%               nearest lattice dot (6.25,4.05) by 0.354, and crosses the
%               middle envelope's TOP edge at x = 6.32, inside that envelope's
%               span 3.29..6.71 -- so it pierces one boundary and not two
%
% MEASURED, not estimated. figs/_harness_clustering.tex \input s this file
% against the real preamble, captures the float into a box with \caption and
% \label gobbled, and \typeout s \wd. Reading:
%   tikzpicture  385.40pt = 135.45mm wide, 187.44pt = 65.87mm tall
%   \textwidth   404.03pt = 142.00mm
%   fills 95.4% of the measure, 18.63pt = 6.55mm of slack.
% The build's overfull-hbox gate is 0; this figure contributes none. The
% fullwidth version it replaces drew 148.6mm against a 174mm allowance -- it
% overran the body measure, which is why it was in a fullwidth in the first
% place, and cost that page its margin column to do it.
%
% TYPE. \scriptsize on every label, matching fig:confound three pages earlier
% and fig:residual-energy in chapter 4. Deliberately not raised: this figure is
% read alongside fig:confound, and a lone larger label set would read as a
% different apparatus rather than as a corrected one.
% ===========================================================================
\begin{figure}[htbp]
\centering
\begin{tikzpicture}[x=1cm, y=1cm, font=\small]
  % --- the constants everything else is derived from -----------------------
  \def\px{1.25}   % column pitch  (one well per column)
  \def\py{0.81}   % row pitch     (one cell line per row)
  \def\hc{4}      % highlighted column index, 0..8
  \def\hr{4}      % highlighted row index,    0..5
  \def\bt{0.25}   % half band thickness
  \def\ep{0.46}   % envelope padding around its three columns
  \def\vm{0.18}   % envelope margin above and below the well band

  \tikzset{
    cond/.style={fill=quietink},
    band/.style={fill=rulegrey!30},
    drugenv/.style={draw=quietink, line width=0.5pt, rounded corners=4pt},
    lead/.style={draw=quietink, line width=0.4pt},
    lbl/.style={font=\scriptsize, text=quietink},
    strong/.style={font=\scriptsize\bfseries, text=accent}}

  % centre lines of the two bands: ON the lattice, by construction
  \pgfmathsetmacro{\HX}{\hc*\px}            %  5.000
  \pgfmathsetmacro{\HY}{\hr*\py}            %  3.240
  % band ends: half a pitch past the first and last dot, both sides, both bands
  \pgfmathsetmacro{\BL}{-\px/2}             % -0.625
  \pgfmathsetmacro{\BR}{8*\px+\px/2}        % 10.625
  \pgfmathsetmacro{\BB}{-\py/2}             % -0.405
  \pgfmathsetmacro{\BT}{5*\py+\py/2}        %  4.455
  \pgfmathsetmacro{\EB}{\BB-\vm}            % -0.585
  \pgfmathsetmacro{\ET}{\BT+\vm}            %  4.635

  % ---- 1. the two crossed groupings, as solid tint bands ------------------
  \fill[band] (\BL,\HY-\bt) rectangle (\BR,\HY+\bt);      % one cell line
  \fill[band] (\HX-\bt,\BB) rectangle (\HX+\bt,\BT);      % one well

  % ---- 2. the drug, as an outline envelope over three whole well-columns --
  % It encloses the well band top and bottom as well as left and right, so the
  % well reads as INSIDE the drug and not merely beside it.
  \foreach \g in {0,1,2}{
    \pgfmathsetmacro{\GL}{3*\g*\px-\ep}
    \pgfmathsetmacro{\GR}{(3*\g+2)*\px+\ep}
    \draw[drugenv] (\GL,\EB) rectangle (\GR,\ET);}

  % ---- 3. the conditions --------------------------------------------------
  \foreach \c in {0,...,8}{
    \foreach \r in {0,...,5}{
      \fill[cond] (\c*\px,\r*\py) circle (1.7pt);}}

  % ---- 4. the intersection ------------------------------------------------
  \fill[accent] (\HX,\HY) circle (2.4pt);
  \draw[accent, line width=0.4pt]
    (\HX+0.11,\HY+0.13) -- (\HX+1.66,\HY+1.75);
  \node[strong, anchor=west] at (\HX+1.78,\HY+1.81)
    {one condition: the intersection};

  % ---- 5. direct labels, one per grouping. No legend, no key. -------------
  % cell line: to the left of its own band.
  \node[lbl, anchor=east] at (\BL-0.15,\HY) {one cell line};

  % well: below its own band, on a lead that PIERCES the envelope's lower
  % edge -- the crossing of that boundary is the nesting, drawn.
  \draw[lead] (\HX,\BB) -- (\HX,\BB-0.34);
  \node[lbl, anchor=north] at (\HX,\BB-0.40) {one well};

  % drug: above the leftmost envelope, centred on it. Placed there and not on
  % the middle envelope because the middle one is where the well band and the
  % callout leader already are, and "one drug" centred over the highlighted
  % column would compete with "one well" for the same column.
  \node[lbl, anchor=south] at (\px,\ET+0.06) {one drug};

  % NO free-floating note line. The version this replaces carried "neither
  % grouping is nested in the other" as a caption-width sentence under the
  % lattice. Two reasons it is gone: the caption now states it in the same
  % words, and at page scale a \scriptsize sentence sitting 4mm above
  % "Figure 3.3." reads as the caption's first line rather than as part of the
  % drawing. The claim is not dropped -- it moved into the caption, where the
  % new drug clause needs it anyway.
\end{tikzpicture}
\caption[The two crossed groupings, and the drug the wells nest in]{\textbf{A
condition belongs to a well and to a cell line at once, and the two groupings
cross.} Each dot is a condition; the horizontal band is one cell line, the
vertical band one well. The dot where they meet is one condition, which is why
the intersection term is the ordinary independent-sampling variance. The
rounded envelopes are drugs: a well carries one drug at one dose, so wells nest
inside drugs, while the cell-line band crosses every drug and neither of the
two bands is nested in the other. The lattice is schematic --- nine columns of
six, for legibility, and not a count of anything in the cache.}
\label{fig:clustering}
\end{figure}


Every interval on a \emph{mean} is therefore two-way cluster-robust
(Cameron--Gelbach--Miller),
\begin{equation}
 V \;=\; V_{\textrm{line}} \;+\; V_{\textrm{well}} \;-\; V_{\textrm{intersection}},
\end{equation}
where the intersection of the two groupings is the single condition, so
$V_{\textrm{intersection}}$ is the ordinary independent-sampling variance. When
it goes non-positive in a finite sample, the wider one-way variance is used and
the branch recorded. The formula above is valid only where the two groupings
genuinely cross, which cell line and well do: a treated well spans
$\approx 49$ cell lines and a cell line recurs across wells. Drug and well do
\emph{not} cross, because one drug is assigned per well and wells therefore nest
strictly inside drugs. Where an interval is reported on the drug axis the
estimator tests the two groupings for nesting first and, finding it, clusters
one-way on the drug rather than applying the expression above to a nested pair;
the branch taken is recorded beside the interval. Reading ``two-way'' as
``the equation above was applied to drug and well'' would therefore be wrong.
It is asymptotic in clusters, so the reference distribution
is Student's $t$ with one fewer degree of freedom than the smaller cluster
count. The \texttt{unseen\_drug} arm shows what that costs: \cref{tab:scale}
puts $199$ scored conditions behind it, drawn from $28$ treated wells and $40$
cell lines, so the smaller cluster count is the $28$ wells, the reference is
$t_{27}$ and not the normal, and the interval widens by $4.7\%$.

For statistics over pairs of conditions, such as the transfer coefficient,
neither grouping partitions the pairs, so those intervals use the
Fafchamps--Gubert dyadic estimator, under which two pairs are dependent when
they share either member.

Throughout, an estimate is quoted as value [lower, upper], a two-way
cluster-robust $95\%$ interval unless stated otherwise. DrEval~\cite{dreval}
identifies the same structure more broadly, citing Hurlbert~\cite{hurlbert}:
hundreds of thousands of cells resolving to a few hundred conditions over $50$
cell lines.

\section{The generation budget is part of the claim}

The second rule concerns generation validity. The budget is a claim-level
variable, not a runtime setting.

A $600$-token cap truncated $26\%$ of predictions in an early run, silently,
because nothing measured them. A prediction cut at the cap never reaches \DOWN{}
and is scored with its whole down-block missing, and no completion rate was
quoted beside the score.

\paragraph{What it cost} The optimal-transport arm of
\cref{sec:res-epsilon}, scored at that cap, gave a gap of
\ci{+0.0144}{-0.0051}{+0.0338} against its geometry-defined \texttt{orth}
partner, whose observed score is near chance. That interval spanned zero and the
arm was printed as inert.\corrmark{The printed negative is retracted in
\cref{sec:res-encoding}, in full and where it was made.} Re-run at the $1400$
tokens the residual arm had, it gives \ci{+0.0292}{+0.0144}{+0.0441}, and the
negative had to be withdrawn (\cref{sec:res-encoding}). The same cap had already
halved the residual arm's own effect. Wherever a paired comparison is scored on
variable-length predictions, then, the token budget is part of the claim, and a
null quoted without its budget cannot be read.

\paragraph{The rule that follows} Every generation-based evaluation computes the
\ENDCELL{}/\DOWN{} completion rate, the fraction of emitted symbols that are
valid panel genes, and the duplicate rate before any score is read from it, and
no score below comes from a run whose validity block was not inspected.

The intervals do not cover the generation configuration. Sampling is at
temperature $0.8$ with top-$p$ $0.9$ throughout, over four budgets: $1200$
tokens with eight predictions per condition in the expression-frame evaluations
feeding \cref{sec:res-blind} and \cref{tab:epsilon}; $1400$ with four in the
residual-frame evaluations of \cref{sec:res-encoding}; $600$ with four in the
control arms until re-run at the residual arm's budget; and $3800$ for output
invariance, which compares whole predictions and is not truncated to a scoring
panel. Where a null turns on its budget the budget is named there, and every
contrast uses a fixed per-(condition, arm) seed.

% PLACEMENT: at the previous setting this box split across the chapter's last
% page break, carrying two lines onto a page that is otherwise almost empty.
% The box is ~9 lines; ask for enough room to keep it whole.
\needspace{10\baselineskip}
\begin{scopelimit}[carried into every residual-frame number below]
Every residual-frame number comes from one generation run and its companion
audits (\texttt{re\_v3} with the scramble-stratum audit, the matched variance
decomposition, the two calibration runs and the channel gate). Numbers from the
two earlier runs are not reconciled with these and are not quoted, except where
a superseded estimate is named as such and the conclusion drawn from it
withdrawn, because between runs the sign of several estimates changed and not
only their magnitude.
\end{scopelimit}

Two components carried defects found after that run and have since been repaired
and re-run, the $\DRF$ calibration bootstrap (\cref{sec:res-metrics}, both arms)
and the channel-gate partner pool (\cref{sec:res-scope}); both are reported on
the corrected runs.

That is the apparatus. It settles what is measured, on what population, against
what reference and with what uncertainty, and reaches no verdict about the
model.
